\scanpage{161}
\inserted{$G$ est une} Extension canonique de $\unsure{\Gamma|\gamma}$ par \inserted{le tore} $G^0$,
\diagram{\begin{tikzcd}
0 \arrow[r,no head] & G^0 \arrow[r,no head] & G \arrow[r,no head] & \Gamma \arrow[r,no head] \arrow[d,equal] & 0 \\
&&& \operatorname{Pic}(A)^{\unsure{\mathrm{abs}|\mathrm{ab}}}
\end{tikzcd}}

Ceci posé, à toute place $\xi$ \deleted{de} $K$, est associé des foncteurs
\diagram{\begin{tikzcd}[column sep=large]
\mathcal M_0 \arrow[r,very thick,"{\text{fibre en }\unsure{\xi}}","{\substack{\sim\\\text{équiv.}}}"'] \arrow[d,equal] \arrow[rr,bend right=42,"{\text{fibre en }\xi}"'] & \substack{\text{Vectoriels sur}\\\mathbf Q\text{ avec}\\\text{opération de }G} & \substack{\text{\unsure{réseaux} de Hodge}\\(+\text{ structure}\\\text{réelle sur}\\\text{\unsure{le tout}, si}\\\text{la place est}\\\text{réelle})} \\
\substack{\text{motifs}\\\text{abéliens}}
\end{tikzcd}}

d'où un foncteur canonique
\diagram{\begin{tikzcd}[column sep=huge]
\substack{\text{Vectoriels sur }\mathbf Q\\\text{avec op. de }G} \arrow[r,"\Phi_\xi"] & \substack{\text{\unsure{réseaux} de Hodge}\\(+\text{ structure réelle}\\\text{si }\unsure{\xi}\ \deleted{\text{\unsure{réel}}}\inserted{\text{réelle}})}
\end{tikzcd}}

\deleted{\unsure{Le} $\Phi_\xi$} Ainsi, si $V$ est un \deleted{\unsure{de}} $\mathbf Q$-vectoriel sur lequel $G$ opère, \uline{alors pour tte place} $\xi$ \uline{de} $K$, \uline{on \unsure{donne}} \uline{une structure de Hodge sur} $V\otimes_{\mathbf Q}\mathbf C$. Il suffit d'ailleurs d'avoir les op. de $G^0$ sur $V$, i.e. de $\uuline{K}^*$ sur $V$ (avec action triviale de $E$), $+$ une place $\widetilde\xi$ de \textcircled{$\widetilde K$}, pour en conclure la \inserted{\uline{structure de}} \unsure{réseau} de Hodge.
\margin{$\widetilde K=$ corps de classes \uline{absolu}}
\margin{160}
